\chapter{Methods}

\section{Ice sheet modelling}
\sectioncount{Ice sheet modelling}

Numerical ice sheet models are important tools for Antarctic science, both for understanding the mechanisms behind recent ice sheet change (e.g.~\citealt{payneRecentDramaticThinning2004,gudmundssonInstantaneousAntarcticIce2019,bradleyQuantifyingAttributingRole2025}), exploring possible ice sheet configurations in the deep past (e.g.~\citealt{deboerSimulatingAntarcticIce2015}) and for making projections of future ice sheet change (\citealt{seroussiISMIP6AntarcticaMultimodel2020,seroussiEvolutionAntarcticIce2024}, among many others). This section reveiws the construction of a general ice sheet model and the important choices involved in simulating the Antarctic ice sheet.

\subsection{Stokes flow}
\sectioncount{Stokes flow}

\noindent
Practical ice flow modelling is done using a hierarchy of approximations to incompressible Stokes flow, which is itself the low Reynolds number limit of the Navier-Stokes equations. In this limit, inertial terms are ignored, and flow is governed by a balance between pressure gradients, viscous stresses, and gravitational driving stresses. These are solved alongside a constitutive relation, which defines the ice rheology (see Section~\ref{sec:rheology}). It is also necessary to impose boundary conditions at the base and surface of the ice sheet. For a rigid boundary (the interface at the bed), this might take the form:

\begin{equation}
    \mathbf{u}_b = F (\boldsymbol{\tau}_b), \quad \mathbf{u}.\mathbf{n} = -m_b,
\end{equation}

\noindent
where $\mathbf{u}$ is the ice velocity field, $\mathbf{u}_b$ is the velocity component parallel to the bed, $\mathbf{n}$ is the normal to the bed, $\boldsymbol{\tau}_b$ is the basal shear stress, $F$ is a sliding law (see Section~\ref{sec:sliding}), and $m_b$ is the basal melt rate (in velocity units). For a free surface (the interface with the atmosphere or ocean), dynamic and kinematic boundary conditions are applied to specify the boundary stresses and surface evolution, respectively. Some model setups include an advection-diffusion equation for temperature, which couples to the momentum (Stokes) equations through the constitutive relation (Eq.~\ref{eq:glenslaw}). However, assuming isothermal ice (i.e.~fixed temperature) is often a reasonable approximation over short timescales.

\subsection{Constitutive relation}\label{sec:rheology}
\sectioncount{Constitutive relation}

The constitutive relation describes the relationship between the stress and deformation of a material. For ice, it is common to use Glen's flow law \citep{glenCreepPolycrystallineIce1955}:
\begin{equation}
    \label{eq:glenslaw}
    \dot{\epsilon}_{ij} = A \tau^{n-1} \tau_{ij}, \quad \tau = \sqrt{\frac{1}{2}\tau_{ij}\tau_{ij}},
\end{equation}

\noindent
where $\epsilon_{ij}$ is the strain rate tensor, $\tau_{ij}$ is the deviatoric stress tensor, $A(T)$ is a temperature-dependant rate factor, and $n$ is a scalar exponent (``Glen's exponent''). Early work to estimate the value of $n$ for polycrystalline ice suggested that $n \approx 3$ (\citealt{weertmanCREEPDEFORMATIONICE1983} and references therein). Subsequently, $n=3$ became standard across a wide range of ice sheet models (e.g.~\citealt{larourContinentalScaleHigh2012,cornfordAdaptiveMeshFinite2013,lipscombDescriptionEvaluationCommunity2019}). However,~\cite{goldsbySuperplasticDeformationIce2001} demonstrated that creep can be divided into multiple regimes, occurring through different mechanisms within the polycrystalline structure. The appropriate value of $n$ depends on the dominant creep mechanism, which in turn depends on the stress.~\note{I could talk here about typical driving stresses in Antarctica and the corresponding value $n=1.8$, but my attempts at this have gotten a bit bogged down and wordy, so have removed for better flow}. Other studies suggest that observations of Antarctic ice streams and ice shelves may be better explained by a higher value of $n$ (e.g.~\citealt{cuffeyHowNonlinearCreep2011,millsteinIceViscosityMore2022}). Ice sheet models that have explored this found Antarctic sea level contribution to be sensitive to the choice of $n$, with higher values leading to greater sea level rise \citep{rosierCalibratedSeaLevel2025,getraerIncreasingGlenNye2025,martinRatesSeaLevelRise2026}.~\note{Is signposting like ``We will explore this uncertainty in chapter 2'' (or similar) appropriate in these situations?}

\todo{Cuffey and Paterson and have a good section on this --- worth a brief rewriting after reading}

\subsection{Approximations in ice flow models}
\sectioncount{Approximations in ice flow models}

Full Stokes modelling is computationally expensive and so is rarely deployed for continent-scale ice sheet modelling (with some notable exceptions, e.g.~\citealt{seddikSimulationsGreenlandIce2012}). Approximations are typically applied to match the requirements of a particular ice-dynamical regime.

\subsection{Basal sliding}\label{sec:sliding}
\subsection{Surface mass balance*}
\subsection{Sub-shelf melting}
\subsection{Glacial isostatic adjustment}
% initial uplift rate

% Most continent-scale projections by coupled ice sheet-sea level models use only a 1D solid Earth model with rheological properties varying only with depth (e.g.~\citealt{gomezSealevelFeedbackLowers2015,larourSlowdownAntarcticMass2019,hanCapturingInteractionsIce2022,hanImprovingProjectionsAntarctic2025}).~\cite{coulonContrastingResponseWest2021} use a modified solid Earth model with a dual pattern for Earth structure under the WAIS and EAIS.\@ However,~\cite{gomezInfluenceRealistic3D2024} have produced the only Antarctic projections accounting for the full 3D Earth structure derived from seismic measurements. They highlight a two-timescale response in far-field sea level

% Both~\cite{coulonContrastingResponseWest2021} and~\cite{gomezInfluenceRealistic3D2024} show that 1D models either underestimate uplift in West Antarctica, overestimate it in East Antarctica, or both.~\cite{gomezInfluenceRealistic3D2024} 

% In continent-wide simulations of Antarctica, the treatment of GRD effects varies across ice sheet models and is an area of active development. Many models do not account for the spatially-varying structure of the solid Earth, instead using only a 1D profile of lithosphere and mantle properties that is spatially homogeneous (e.g.~\citealt{gomezSealevelFeedbackLowers2015,larourSlowdownAntarcticMass2019,hanCapturingInteractionsIce2022,hanImprovingProjectionsAntarctic2025}).

\subsection{Initialisation}

